% =====================================================================
% DRAFT: Positioning section for QVCache vs. result-materialized caches
% (Aker, Proximity, Potluck, GPTCache).
%
% Suggested placement: new Section 3 ("The Design Space of Approximate
% Caches for Vector Search"), immediately before the current Solution
% Overview, with the short related-work paragraph at the end folded into
% Section 7.
%
% Claims marked [VERIFY] are deductions from published algorithms and
% from the Aker reference implementation; they must be backed by a
% measurement before submission. See notes at the end of the file.
% =====================================================================

\section{The Design Space of Approximate Caches for Vector Search}
\label{sec:design-space}

Prior work distinguishes \emph{traversal-centric} caches, which retain
index structures and vectors touched during ANN traversal, from
\emph{result-centric} caches, which retain the outcomes of past
queries~\cite{aker}. This distinction captures \emph{what} is stored,
but it obscures a second choice that we argue is more consequential:
given that a cache stores the outcomes of past queries, does it store
those outcomes as \emph{materialized answers} or as \emph{materialized
data}?

We therefore refine the taxonomy. Let $P \subset \mathbb{R}^d$ be the
indexed collection and let $q_1,\dots,q_N$ be the queries a cache has
observed, with $A_i = \mathrm{kNN}_P(q_i)$ the top-$k$ each returned.

\begin{itemize}
  \item A \textbf{result-materialized} cache stores the pairs
        $(q_i, A_i)$ and answers a new query $q$ by \emph{selecting}
        one $A_i$ for reuse. Aker~\cite{aker}, Proximity~\cite{proximity},
        Potluck~\cite{potluck}, and GPTCache~\cite{gptcache} are all of
        this form.
  \item A \textbf{data-materialized} cache stores the union
        $S = \bigcup_i A_i \subseteq P$ in a searchable index and answers
        $q$ by \emph{recomputing} $\mathrm{kNN}_S(q)$. QVCache is of this
        form.
\end{itemize}

Both designs store the same information: after the same query sequence,
the vector content of $S$ and of a result-materialized cache's result
pool are identical up to eviction policy. They differ only in the
\emph{access path} they provide over that content --- selection versus
search. This section shows that this single difference determines the
reuse geometry a cache can exploit (\S\ref{sec:ds-geometry}), the
information its hit predicate can observe (\S\ref{sec:ds-predicate}),
and the cost of maintaining it under updates
(\S\ref{sec:ds-updates}). \S\ref{sec:ds-limits} states the regimes in
which result-materialization remains preferable.

% ---------------------------------------------------------------------
\subsection{Reuse Geometry: Atomic versus Compositional}
\label{sec:ds-geometry}

Write $r_j(q)$ for the distance from $q$ to its $j$-th nearest neighbor
in $P$, so that the exact answer to $q$ is contained in the ball
$B(q, r_k(q))$.

\paragraph{Result-materialized reuse is atomic.}
A result-materialized cache admits $q$ when $q$ lies within a reuse
radius $\tau_i$ of some cached query, and then returns $A_i$ unchanged.
Its servable region in query space is therefore a union of balls
centered on \emph{previously issued queries},
\begin{equation}
  \mathcal{C}_{\mathrm{res}} \;=\; \bigcup_{i=1}^{N} B(q_i, \tau_i),
  \label{eq:cov-res}
\end{equation}
and its answer space is the finite set $\{A_1,\dots,A_N\}$: over the
lifetime of the cache, at most $N$ distinct responses are reachable.
The radii in \eqref{eq:cov-res} are necessarily small. Aker initializes
$\tau_q$ to a quarter of the distance from $q$ to the nearest member of
its own cached result and caps it above by that distance, so
$\tau_i = \Theta(r_1(q_i))$.

\paragraph{Data-materialized reuse is compositional.}
A data-materialized cache evaluates $\mathrm{kNN}_S(q)$ at query time.
Its answer space is every $k$-subset of $S$, and its servable region is
the set of queries for which $S$ is locally an adequate surrogate for
$P$,
\begin{equation}
  \mathcal{C}_{\mathrm{data}} \;\approx\;
  \bigl\{\, q \;:\; r^{S}_k(q) \le (1{+}D)\,\theta[k][R(q)] \,\bigr\},
  \label{eq:cov-data}
\end{equation}
which, because $S \supseteq \bigcup_i \bigl(P \cap B(q_i, r_k(q_i))\bigr)$,
is approximately a neighborhood of the union of past \emph{result}
balls rather than of past query points.

\paragraph{Two consequences.}
First, the governing radius differs: \eqref{eq:cov-res} is built from
balls of radius $\Theta(r_1)$ and \eqref{eq:cov-data} from balls of
radius $r_k$. Because distances concentrate in high dimensions, the
ratio $r_k/r_1$ is modest, and combined with the constant-factor
conservatism above the radius ratio is a single-digit number. Coverage,
however, is a volume, so a single-digit radius ratio corresponds to a
ratio of order $c^d$ in covered space. At $d \ge 100$ the union in
\eqref{eq:cov-res} is, for practical purposes, a set of isolated
points.

Second, and more fundamentally, the balls in \eqref{eq:cov-data}
\emph{compose} and those in \eqref{eq:cov-res} do not. Consider two
cached queries separated by $3\tau$ and a new query $q$ at their
midpoint. A result-materialized cache misses, because $q$ lies outside
both balls --- yet $A_1 \cup A_2$ may contain the exact
$\mathrm{kNN}_P(q)$, and those vectors are \emph{already resident in the
cache}. The limitation is not what the cache stores but that its access
path exposes one frozen list at a time. A data-materialized cache
reaches the same vectors by construction, and can return a $k$-set
assembled from the contributions of many past misses --- a set that no
single previous query ever returned.

This yields two testable predictions. Because the per-miss covered ball
has radius $r_k$ for data-materialized caches and $\Theta(r_1)$ for
result-materialized ones, the advantage should \emph{grow with $k$};
this is consistent with Figure~\ref{fig:varying-k}, where QVCache's
latency reduction increases with $k$. And because compositional reuse
requires only that the workload's queries draw their neighbors from a
shared region of $P$, not that individual queries recur, it should
persist where pairwise query similarity vanishes. Section~\ref{sec:noise}
confirms this: at a perturbation ratio of $\eta = 0.6$, where perturbed
queries share \emph{no} top-$k$ neighbors pairwise, QVCache still
sustains a hit ratio near $0.75$. Under \eqref{eq:cov-res} this regime
admits essentially no approximate hits, since
$d(q,q') \gtrsim r_k > r_1 > \tau$ for nearly every pair.

% ---------------------------------------------------------------------
\subsection{What the Hit Predicate Can Observe}
\label{sec:ds-predicate}

A cache must decide whether a candidate answer is good enough before
returning it. The two designs differ in what evidence is available at
that moment.

Because a result-materialized cache holds a frozen list, it cannot
assess the quality of the answer it is about to return with respect to
the incoming query; it can only measure how far that query has moved
from the one that produced the list. Its predicate is thus a test on a
\emph{proxy}, $d(q,q_i) \le \tau_i$. A data-materialized cache computes
$\mathrm{kNN}_S(q)$ first and tests the realized quantity,
$r^{S}_k(q)$, directly against a learned estimate of $r^{P}_k(q)$. Query
displacement is a conservative sufficient condition for answer quality,
but the converse fails routinely: many queries far from every $q_i$
nonetheless admit excellent cached answers. Consequently
\eqref{eq:cov-res} is an inner approximation of the true adequacy
region, and it is an inner approximation of the wrong shape, since the
adequacy region is not centered on past queries at all.

\paragraph{The adaptation signal.}
This asymmetry propagates into how each design learns its threshold.
QVCache updates $\theta[k][R]$ from $d_{\mathrm{backend}}[k]$, a
measurement of the quantity it is trying to estimate, obtained from the
backend on every cache miss (Algorithm~\ref{alg:learn}). The signal is
\emph{countercyclical}: it is most abundant precisely when the cache is
least effective, so the control loop cannot starve, and it requires only
that queries revisit the same region --- not that any query recur.

Aker instead adapts $\tau_q$ from the \emph{hit type}: it multiplies
$\tau_q$ by $\alpha_L = 1.1$ on an equality hit and by
$\alpha_T = 0.25$ on an approximate hit~\cite{aker}. Hit type is not a
measurement of accuracy; no ground-truth comparison is performed on
either branch. (Potluck's dropout mechanism~\cite{potluck}, by contrast,
does verify against the backend, at the cost of forced misses.) Two
properties follow. Sustaining a given $\tau_q$ requires
$\ln(1/\alpha_T)/\ln(\alpha_L) \approx 14.5$ equality hits per
approximate hit. And in the reference implementation, only a
\emph{bit-exact} repeat of a representative query loosens the
threshold; a repeat of a query previously served approximately is
resolved through an alias entry and contributes no loosening
signal.\,[VERIFY]

In workloads whose query vocabulary is effectively unbounded --- the
regime that motivates similarity caching in the first place, since
distinct user inputs yield distinct embeddings
(\S\ref{sec:workload-characteristics}) --- equality hits are rare, and
$\tau_q$ therefore decays monotonically: from $\texttt{minD}/4$ to
$\texttt{minD}/16$ after one approximate hit and to $\texttt{minD}/64$
after two, with no lower bound. Adaptation then degenerates into a
monotonically shrinking fraction of the local nearest-neighbor
distance, and the density-awareness of the mechanism reduces to the
local scaling already present in its initialization.

\paragraph{On evaluation regimes.}
This effect is invisible under workloads with a small finite query
vocabulary. Aker's \textsc{simZipf} generator draws each request
uniformly from a pre-generated group of $50$ variants per base query, so
a frequently accessed group repeats each of its members many times, and
the authors report that equality hits dominate in their
evaluation~\cite{aker}. Because equality hits are by construction free
of reuse error, this regime is favorable to hit-type adaptation on both
the hit-ratio and the recall axis. We therefore evaluate the
complementary regime as well, treating vocabulary size as an explicit
workload parameter (\S\ref{sec:vocab-sweep}): our generator
synthesizes fresh perturbed queries rather than resampling a fixed pool,
so no two requests are identical.

% ---------------------------------------------------------------------
\subsection{Maintenance Under Updates}
\label{sec:ds-updates}

The two designs also diverge sharply once the underlying collection
changes, and again the divergence is a consequence of the access path
rather than of any particular refresh algorithm.

\paragraph{Refresh cost is a product versus a sum.}
In a result-materialized cache each frozen list must be repaired
independently, so propagating $M$ insertions to $N$ cached queries is
$O(MN)$ distance evaluations in the worst case. Aker manages this cost
by deferring most of it: a fast path repairs the single entry nearest to
the inserted vector, and a slow path advances one entry at a time
through a batch of $B$ insert-log records, gated on a global risk
score~\cite{aker}. Full propagation of one insertion therefore requires
$O(N)$ slow-path invocations, is conditional on the risk score
exceeding its threshold, and pins insert-log memory behind whichever
entry lags furthest. In a data-materialized cache an insertion is a
single index insertion; total work is $O(M)$, independent of how many
queries the cache has served.

\paragraph{Visibility scope.}
The scopes also differ. Repairing a frozen list makes the new vector
available only to queries near that one cached query; inserting into $S$
makes it available to \emph{every} subsequent query, including queries
with no similar predecessor in the cache. We report this as \emph{insert
visibility lag}: the interval between a committed backend insertion and
the point at which the cache would return the new vector to a query
whose exact top-$k$ contains it.

\paragraph{Deletions.}
Under deletions, a result-materialized cache faces a structural cliff:
removing one vector shortens every list that references it, and any list
left with fewer than $k$ valid members must be invalidated wholesale.
Aker mitigates this by storing $\Delta$ reserve neighbors per entry,
which consumes pool budget and so reduces the number of cached entries
--- a trade-off it cannot escape. A data-materialized cache has no
entries to invalidate. Deleting a cached vector removes one node from a
graph over $|S|$ nodes, and since graph search already maintains a
candidate pool of size $L \ge k$, the top-$k$ is backfilled from the
surviving cached vectors at no additional memory cost; the effective
reserve is $L-k$ rather than a provisioned $\Delta$. Degradation is
proportional to the fraction of cached vectors deleted rather than
discontinuous in it.

\paragraph{A stronger consistency model.}
Aker's \emph{del-consistency} pairs validity ($c_t \subseteq V_t$,
$|c_t| = k$) with \emph{relaxed visibility}, which places no bound on
insertion staleness~\cite{aker}. We argue that relaxed visibility is not
a design preference but the strongest guarantee a result-materialized
cache can afford, since bounding insertion staleness there requires
paying the $O(MN)$ product. Data-materialization permits a strictly
stronger model, which we call \emph{live-consistency}:

\begin{description}
  \item[L1 (Validity).] At any time $t$ a cache hit returns only live
    vectors, and returns $k$ of them: $c_t \subseteq V_t$,
    $|c_t| = k$.
  \item[L2 (Region freshness).] If $v$ is inserted at $t_0$ into a
    region that is active at $t_0$, then within a bounded admission
    delay $\delta$ --- one index insertion --- $v$ is visible to
    \emph{all} queries in that region for as long as it remains
    resident.
  \item[L3 (Fail-to-miss).] Whenever freshness cannot be established
    for a region, queries in that region degrade to a cache
    \emph{miss}, never to a stale hit.
\end{description}

L3 is what decouples accuracy from update rate. Under
del-consistency, refresh lag manifests as a stale answer, so recall is a
function of the update rate; under live-consistency it manifests as a
backend query, so only the hit ratio is. Moreover, given L1 and L2, a
cache hit returns exactly the top-$k$ that the backend would return over
the live set restricted to $S$ --- and therefore the recall analysis of
\S\ref{sec:hit-miss}, which concerns the quality of $\theta$ as an
estimator of $r^P_k$, carries over unchanged to the dynamic setting
\emph{once L3 holds}. The write-path rule below is how L3 is
implemented; it is not a second hit test.

\paragraph{Realizing L1--L3.}
Write $V_t$ for the live collection at time $t$ and
$S_t\subseteq V_t$ for the live vectors resident in the mini-indexes
(tombstones already excluded). Let \(R:\mathbb{R}^d\to\{0,\ldots,B{-}1\}^{d'}\)
be the regional key (PCA to \(d'{=}16\), \(B{=}8\) bins) and let
\(\theta_t[R]\) be the learned estimate of \(r^{V}_k\) for queries with
that key. A query \(q\) is a hit iff \(S_t\) returns \(k\) live ids and
\(r^{S_t}_k(q)\le(1{+}D)\,\theta_t[R(q)]\). A hit is
\emph{stale} when
\(\mathrm{kNN}_{S_t}(q)\ne\mathrm{kNN}_{V_t}(q)\). L3 is exactly the
requirement that stale hits do not occur.

Axiom L1 is eager: a delete of \(v\) tombstones \(v\) in every shard
before the next search, so a hit cannot return a dead id. Axiom L2 is
write-path admission when \(S\) has a free slot --- one index insertion,
then compositional visibility: any later \(\mathrm{kNN}_{S}(q)\) may
include \(v\), not only queries near the writer. Writes never evict a
shard to make room; if \(S\) is full, L2 is deferred to the next miss in
that region (miss-admit), rather than destroying unrelated
neighborhoods.

Axiom L3 cannot inspect \(\mathrm{kNN}_{V_t}(q)\) at write time. A
write of \(v\) at \(t_0\) can stale only queries whose live top-$k$
contains \(v\) (insert) or contained \(v\) (delete). Call that
\emph{influence set}
\begin{equation}
  \mathcal{I}(v)
  \;=\;
  \bigl\{\, R(q) \;:\;
    v\in\mathrm{kNN}_{V_{t_0}}(q)
    \;\text{or}\;
    v\in\mathrm{kNN}_{V_{t_0^-}}(q)
  \,\bigr\}.
  \label{eq:influence}
\end{equation}
\(\mathcal{I}(v)\) is not observable without a query log. The
implementation replaces it by a locality cover on the region code. Let
\(B_1(R(v),r)\) be the \(\ell_1\) ball of Manhattan radius \(r\) on the
\(d'\)-byte bucket key (the write radius; not to be confused with
axiom~L1). The working hypothesis is that \(k\)NN relationships are
locally Lipschitz in this code: there exists a small \(r_\star\) such
that \(\mathcal{I}(v)\subseteq B_1(R(v),r_\star)\) for writes that
matter. Then, for every \emph{already learned} cell
\(R\in B_1(R(v),r)\),
\begin{equation}
  \theta[R] \;\leftarrow\; \alpha\cdot\theta[R],
  \qquad \alpha\in(0,1),
  \label{eq:theta-penalty}
\end{equation}
once, until a miss EMA-updates that cell (a saturating freshness bit).
Search-only never applies \eqref{eq:theta-penalty}; the hit formula
itself is unchanged.

Why this is fail-to-miss, not a new predicate. An insert of a closer
neighbor \emph{decreases} \(r^{V}_k\) while \(r^{S}_k\) stays at the
pre-write value if \(v\notin S\) or is poorly connected. The old
\(\theta\) still licenses \(r^{S}_k\le\theta\), which is a stale hit.
Contracting \(\theta\) shrinks the servable set
\eqref{eq:cov-data} in every influenced cell, so those queries miss,
the backend observes the live \(d_k\), and the EMA restores a fresh
estimate. A delete is dual: L1 already hides \(v\); the hole is
backfilled from the surviving graph (\(L-k\) implicit reserve), which
can still satisfy the old \(\theta\) with the wrong \(k\)-set.
\eqref{eq:theta-penalty} removes that slack.

Two parameters govern the cover, not the model. Radius \(r\) is the
completeness of \(\widehat{\mathcal{I}}(v)=B_1(R(v),r)\) as a
substitute for \eqref{eq:influence}: \(r{=}0\) (exact cell of \(v\))
under-covers, so L3 fails and hit recall collapses after writes;
large \(r\) over-covers and L3 fires on cells that were never
influenced (a miss tax, not a recall bug). The factor \(\alpha\) is
how hard the fail-to-miss is: \(\alpha{=}1\) is a no-op;
\(\alpha\to 0\) is erasure of \(\theta\) (the prefix wipe). Saturating
\(\alpha\) is required by the model: L3 asks for one loss of
freshness, not \(\theta\to 0\). After a miss, \(\theta\) is again an
estimate of live \(r^{V}_k\) and the cell is trusted until the next
write.

Under a cover with \(r\ge r_\star\) and \(\alpha\) small enough that
\(\alpha\,\theta[R] < r^{S}_k(q)\) for every stale \(q\) in those
cells, \eqref{eq:theta-penalty} implements L3 and the recall argument
is independent of the update rate, as claimed above. Finite \(r\) and
\(\alpha{=}0.8\) are a tunable inner approximation of that cover, not
a second consistency model. Empirically, \(r{=}2\) on the 16-D code
recovers hit recall \(\approx 0.91\) after a 5\% insert / 5\% delete
refresh, versus \(\approx 0.80\) at \(r{=}0\), while penalizing
\(\approx 19\)k of \(60\)k learned cells rather than the whole
map. The remaining error is uncovered influence, not a hole in L1.

% ---------------------------------------------------------------------
\subsection{When Result-Materialization Is Preferable}
\label{sec:ds-limits}

The comparison is not uniformly favorable, and three regimes favor
result-materialization.

\emph{Lookup cost.} Selection is $O(1)$: an equality hit is a hash
probe. Composition requires a graph traversal, $O(\log|S|)$. Under
workloads dominated by exact query repetition, a result-materialized
cache attains lower per-hit latency, and its query lookup table and
alias entries are precisely the right design for that regime. QVCache
buys a larger servable region at the price of a more expensive lookup,
a trade that is favorable when backend search costs milliseconds and
unfavorable when it does not.

\emph{Memory per cached vector.} A data-materialized cache stores a
proximity graph over $S$, whose edges cost several hundred bytes per
vector, whereas a result-materialized cache stores raw vectors with
small per-entry metadata and indexes only cache \emph{entries}. The net
comparison depends on $d$, $k$, $\Delta$, and the degree of result
sharing across entries, and is not favorable \emph{a priori}; we
measure it in \S\ref{sec:memory}.

\emph{Estimator quality.} Compositional reuse rests on $\theta[k][R]$
being an adequate estimate of $r^P_k$ throughout region $R$. Where a
region straddles both dense and sparse structure, $\theta$ is
overestimated on the dense side and hits are admitted whose cached
$k$-set is not the exact one. A result-materialized cache's
conservatism instead errs toward misses. QVCache therefore exchanges a
small, safe servable region for a near-maximal, estimator-dependent
one, and the granularity parameters $n_{\text{buckets}}$ and
$d_{\text{reduced}}$ govern that exchange
(\S\ref{sec:sensitivity-partitioning}).

Finally, both designs bound visibility by residency: eviction can
remove a vector that was previously visible. L2 is accordingly stated
as holding while $v$ remains resident.

% =====================================================================
% Short paragraph to fold into Section 7 (Related Work).
% =====================================================================

\paragraph{Approximate caches for vector search.}
Proximity~\cite{proximity} reuses cached results under a fixed global
similarity threshold; Potluck~\cite{potluck} tunes a global threshold at
runtime by probabilistically forcing misses and comparing against the
engine's output; GPTCache~\cite{gptcache} applies the same idea to LLM
prompts. Aker~\cite{aker} advances this line in two directions,
replacing the global threshold with a per-cached-query radius adapted
from observed hit types, and introducing del-consistency together with a
refresh pipeline that applies deletions eagerly and insertions lazily.
QVCache is complementary in mechanism but differs in kind: all of these
systems materialize \emph{answers} and reuse them by selection, whereas
QVCache materializes the retrieved \emph{data} and reuses it by search.
As \S\ref{sec:design-space} develops, this is what lets QVCache compose
a top-$k$ across many past misses, supervise its threshold with backend
distances rather than hit types, and reduce cache refresh to index
maintenance.

% =====================================================================
% NOTES BEFORE SUBMISSION
%
% [VERIFY] 1. The claim that only bit-exact repeats of a representative
%   query loosen tau_q is read from Aker's reference implementation
%   (src/core/ak_anns_cache_similarity_engine.cc: the alpha_loosen update
%   is guarded by !linked_hit; the alpha_tighten update fires on the
%   approximate-filter path). Confirm against the released version you
%   cite, and prefer citing observed behavior over source lines.
%
% [VERIFY] 2. The tau_q decay dynamics are deduced from the published
%   update rule and defaults, not measured. Instrument tau_q over a
%   unique-query workload and report its distribution over time; a
%   histogram collapsing toward zero makes the argument unarguable and
%   removes any reliance on reading their code.
%
% [MEASURE] 3. Sec. ds-limits promises memory numbers. Do not claim a
%   memory win before measuring it -- graph edges plausibly make
%   QVCache more expensive per cached vector.
%
% [FIGURE MAP]
%   Sec. ds-geometry   -> existing varying-k figure; existing eta sweep.
%   Sec. ds-predicate  -> NEW vocabulary-size sweep; NEW tau_q histogram.
%   Sec. ds-updates    -> NEW hit ratio + recall vs deletion rate;
%                         NEW insert visibility lag; NEW per-update
%                         cache-side overhead; NEW threads vs QPS
%                         under a concurrent update stream.
% =====================================================================
